Fix real $\varepsilon>0$ and an integer $k\ge 2$. For $n\in\mathbb{N}$, call an integer $m\ge 1$ \emph{$n^{\varepsilon}$-smooth} if every prime divisor $p\mid m$ satisfies $p\le n^{\varepsilon}$ (equivalently $P(m)\le n^{\varepsilon}$, where $P(m)$ denotes the largest prime factor of $m$ and we set $P(1)=1$).

\medskip
\noindent\textbf{Claim (literal statement).} There exists $n_0(\varepsilon,k)$ such that for every $n\ge n_0(\varepsilon,k)$ there is a block of $k$ consecutive integers contained in $\{1,2,\dots,n\}$, all of which are $n^{\varepsilon}$-smooth.

\medskip
\noindent\textbf{Misstatement/ambiguity note.} The literal claim is elementary because the block may be chosen near $1$ and thus does not ``scale with $n$.'' A \emph{minimal nontrivial correction} (consistent with how such questions are usually posed) is for example:

\smallskip
\noindent\textbf{Corrected variant (one natural choice).} Fix $\varepsilon>0$ and $k\ge 2$. Is it true that for all sufficiently large integers $m$ there exist $k$ consecutive integers $m,m+1,\dots,m+k-1$ such that each is \emph{$m^{\varepsilon}$-smooth} (i.e. $P(m+j)\le m^{\varepsilon}$ for $0\le j\le k-1$)?
